\hypertarget{ux6458ux8981}{%
\chapter{摘要}\label{ux6458ux8981}}

本作业从第7讲课时1第8页ReAct``想---做---看---再想''出发，研究多Agent何时需要重新协调，以及修改自己的安排时如何考虑队友的后续任务。借鉴MPC有限前瞻、近期执行与反馈更新思想，提出各自规划、相关方响应、共同确认的滚动协作设计，并实现同构Agent核心。早期Flatland多响应实验在开发集有改善，但保留集未体现相对单响应的额外任务收益，计算和通信开销较高；Qwen小样本也未证明模型增益。随后缩小验证目标，在真实Flatland中比较保留有效许可、按需协调与周期重复询问。三个已有七车场景、三档扰动、四方法共36局全部到达。按需协调与每步询问在九组逐步轨迹完全相同的条件下，含状态发布的通信字节减少96.53\%，本地响应评估减少97.36\%；相对每8步询问，累计完成步数减少15.06\%。据此，本作业在当前开发条件下验证了保留有效安排、按需协调的降本价值，但没有证明复杂连锁修复、多响应或整个框架的普遍优势，也没有把已有事件触发思想本身视为独有创新。

\hypertarget{ux8bfeux7a0bux5207ux5165ux70b9ux4e0eux95eeux9898}{%
\chapter{课程切入点与问题}\label{ux8bfeux7a0bux5207ux5165ux70b9ux4e0eux95eeux9898}}

课件《第7讲\_课时1\_质量良好是怎么来的》第8页介绍``ReAct节拍''：想一步、做一步、看一步，再根据观察决定下一步。课件以南海北部H3编码及质检查询为例；本作业将该循环延伸到协作计划修订，铁路是验证载体，不声称得到海洋业务效果。

具体问题：A准备改变通行安排，B尚未进入相关区段，但后续路线已经依赖A的让行。仅判断``B是否已经行动''会漏掉未来需求；仅让A预测B又可能误判。因此由B自己回答：这个候选是否可接受，能否改变自己的路线或时间窗，代价多少？C也按相同方式参与。问题不要求出现振荡，关注可行性、延误及修订成本。

\hypertarget{ux76f8ux5173ux601dux60f3ux4e0eux539fux521bux8bbeux8ba1ux8fb9ux754c}{%
\chapter{相关思想与原创设计边界}\label{ux76f8ux5173ux601dux60f3ux4e0eux539fux521bux8bbeux8ba1ux8fb9ux754c}}

ReAct把推理、行动和观察组织为循环；MPC强调有限时域前瞻、只执行近期控制、根据反馈滚动更新。分布式协调和计划修复已有大量研究，本作业不将``消息确认''或``局部修复''单独宣称为新发明。

本次原创探索是把这些思想具体组合为Agent的统一能力：\textbf{候选变更---队友自行响应---公开条件协调---近期确认执行---反馈修订}。原创性指本作业提出和实现的改进设计，非经系统查重确认的全球首次。

与早期思路相比，设计经历三次转变：从只保护旧承诺转向比较新的联合安排；从A预测B转向B自行报价；从固定中央规划转向同构Agent及临时提案角色。真实过程中的用户质疑和修订，是本方案形成的重要部分。

\hypertarget{ux65b9ux6cd5ux8be6ux7ec6ux62c6ux89e3}{%
\chapter{方法详细拆解}\label{ux65b9ux6cd5ux8be6ux7ec6ux62c6ux89e3}}

\hypertarget{ux72b6ux6001ux5bf9ux8c61ux4e0eux89d2ux8272}{%
\section{状态、对象与角色}\label{ux72b6ux6001ux5bf9ux8c61ux4e0eux89d2ux8272}}

每个Agent具有自己的观测、私有计划、承诺版本及上下文，运行同一\passthrough{\lstinline!Agent!}类。公开候选包含通行资源格、进入/退出时刻、有效期和版本。Response包含可行性、增量代价、公开条件、私有响应计划；只有前三者进入消息。

提案方临时汇总公开响应，不获取队友私有路线对象或替它生成计划。当前组合选择有临时协调角色，不称完全无协调者算法。单进程模拟多个Agent是工程实现，不等同真实网络部署；任务适配器持有共享环境引用，但只读取自身列车和公开轨道拓扑，不具有进程级权限隔离。

\hypertarget{ux516dux6a21ux5757ux53caux5b9eux9645ux5b9eux73b0}{%
\section{六模块及实际实现}\label{ux516dux6a21ux5757ux53caux5b9eux9645ux5b9eux73b0}}

\begin{longtable}[]{@{}lll@{}}
\toprule
模块 & 实现方式 & 不能夸大的部分 \\
\midrule
\endhead
承诺期限 & 8环境步有效期，附版本及成立条件 & 不是已学会自适应期限 \\
当前与未来依赖 & 各方比较候选占用及自身32步未来路线，报告通行条件 &
主要是空间时间可行性，未独立证明反事实任务依赖价值 \\
候选前瞻 & 最多3条路线，0/3/8/16步等待，保持项可用时纳入 & 不是完备规划 \\
消息确认 & 指定提案、参与集合、候选与响应编号，准备后提交 &
同步可靠通信，不是异步容错共识 \\
局部修复 & 队友返回最多3项本地响应，安装被确认项 &
单位计划变更代理成本，不等于现实经济损失 \\
紧急撤销 & 故障或观测轨迹偏离则撤销，停止并通知 & 无全局无死锁或安全证明 \\
\bottomrule
\end{longtable}

\hypertarget{ux54cdux5e94ux7ec4ux5408ux4e0eux8fd1ux671fux6267ux884c}{%
\section{响应组合与近期执行}\label{ux54cdux5e94ux7ec4ux5408ux4e0eux8fd1ux671fux6267ux884c}}

\begin{lstlisting}
本地观察 → 是否值得改计划？
          → 公布有限候选
          → 队友各自生成条件化响应菜单
          → 检查公开条件、比较增量总代价
             冲突 → 继续尝试替代组合
          → 各自版本核验 → 提交选中响应
          → 每车执行一个动作 → 新观察
\end{lstlisting}

每候选以代价优先检查最多128个响应组合。发起方仅选择各方已公布的选项，不替其规划；有界搜索可能漏解。这里是一次菜单交换后的组合协调，不是新约束反馈下的多轮局部重新求解。预测时域32步、执行跨度1步、承诺期限8步三者不同。

增量代价以自己的保持计划为基准，计算剩余路线行程与等待变化，并按设置加入每次改变已有预测计划1单位修复成本。不同列车代价求和，硬冲突不能被低代价抵消。通信尚未进入优化目标，单独报告。

\hypertarget{ux53efux63d2ux62d4ux5b9eux73b0}{%
\section{可插拔实现}\label{ux53efux63d2ux62d4ux5b9eux73b0}}

核心循环为观测决策、消息交换、近期行动。\passthrough{\lstinline!Task!}实现observe、candidates、respond/response\_options、compatible、valid/install、revoke、act等接口；\passthrough{\lstinline!Transport!}和\passthrough{\lstinline!Clock!}只负责投递与同步推进。铁路逻辑位于实验目录的\passthrough{\lstinline!flatland\_task.py!}，不写入通用核心。

组内通信和跨组依赖扩展已具备协议实现。早期12局使用七列车单交通组；后续resource-scoped-response-v2按候选及公开未来路线的资源交集选择参与者，可能仍包括全部七名成员。分组收益尚未成立，不将其写作主实验贡献。

\hypertarget{ux6700ux65b0ux65b9ux6848ux4e0eux672cux6b21ux5b9eux73b0ux9010ux9879ux5bf9ux9f50}{%
\section{最新方案与本次实现逐项对齐}\label{ux6700ux65b0ux65b9ux6848ux4e0eux672cux6b21ux5b9eux73b0ux9010ux9879ux5bf9ux9f50}}

本节表格主要描述004历史实验版本，不能作为最新源码能力清单。之后核心增加分段补丁、依赖管理及可选缓存等机制；007验证通用局部连锁，008为铁路离线预约诊断，均不等于在线收益。009另用全路线许可适配器，仅验证保留有效安排及按需协调，未启用完整六模块。不同版本的能力、运行及结论分别引用，不用新代码替旧实验补证。

本表区分``希望怎样做''``代码实际做了什么''和``实验支持到哪一步''。方案更新不等于历史实验已运行该能力。

\begin{longtable}[]{@{}
  >{\raggedright\arraybackslash}p{(\columnwidth - 4\tabcolsep) * \real{0.33}}
  >{\raggedright\arraybackslash}p{(\columnwidth - 4\tabcolsep) * \real{0.33}}
  >{\raggedright\arraybackslash}p{(\columnwidth - 4\tabcolsep) * \real{0.33}}@{}}
\toprule
方案要求 & 当前实现 & 本次证据与剩余边界 \\
\midrule
\endhead
只详细规划近期，执行到需要真实反馈的位置 &
固定预测32步、承诺有效期8步，每次推进1步后重新观察 &
已实现固定步长滚动；尚未按工具结果、故障或任务阶段自动选择反馈边界 \\
A不替B规划，由B判断自己怎样配合 &
每个同构Agent通过respond/response\_options生成本地响应，发起方组合公开条件 &
已有真实提交案例；规则响应为主，不是各Agent持续进行LLM推理 \\
考虑队友尚未行动的未来依赖 & 比较32步预测占用；v2按候选和公开未来路线资源交集纳入成员
& 可覆盖部分未来相关成员，但不等于完整任务因果依赖，也未单独证实该模块收益 \\
没有相互影响就各自执行、不通信 &
v2通过proposal\_participants筛选参与者，使用保守资源交集 &
尚无精确的无外部影响判定及零通信保证；路线共享不一定意味着时间窗冲突，可能多通知 \\
只协调真正共享约束的成员，必要时补齐依赖 &
协议支持参与者扩展；Flatland适配提供资源相关成员 &
尚未验证复杂任务依赖完整性；不能将目的地分组失败或保守交集当作精确分组成功 \\
约束冲突时继续协商 & 每候选最多检查128个现有响应组合，每方最多3项响应 &
已实现有限菜单组合，不支持把新约束反复发回各方并重新生成响应，可能漏解 \\
复用未受影响的计划与评估 & 旧计划可作为保持候选，确认后安装各自计划 &
不代表已有跨轮评估缓存或求解器热启动；未通过专门消融证明重算减少 \\
消息只发送协商所需变化 & 交换候选占用、条件、编号和版本 &
尚无差量编码和接收端版本缓存；现有完整报文较大，不能宣称已降低传输成本 \\
承诺失效可撤销，修改时计入修复代价 &
故障或轨迹偏离触发撤销；已有预测改变计单位代理成本 &
已记录撤销及计划变更，不等于物理恢复完成；保留集加修复代价未改善结果 \\
\bottomrule
\end{longtable}

对应代码为\href{../expr/../核心源码/agent.py}{Agent核心}中的参与者选择与消息循环，以及\href{../expr/004-distributed-agent-core/src/flatland_task.py}{Flatland适配器}中的候选、响应与资源筛选。同步可靠通信、近似轨迹预测是当前前提，见\href{../expr/前提.md}{前提}。因此本文评价的是现有有限菜单实现，而不是对全部设计目标的验证。

\hypertarget{ux5b9eux9a8cux8bbeux7f6eux4e0eux516cux5e73ux6027}{%
\chapter{实验设置与公平性}\label{ux5b9eux9a8cux8bbeux7f6eux4e0eux516cux5e73ux6027}}

环境为隔离Python3.11.15、Flatland4.2.6、NumPy1.26.4。数据由官方场景仓库固定commit的Test\_00元数据配置生成，v1含Level\_0/1/2，每例7列车、不同速度及原生故障；不是官方排行榜原始测试集。来源commit与SHA256见\href{../expr/datasets/flatland-core/v1/manifest.json}{datasets
manifest}。

证据分三层。早期\passthrough{\lstinline!response-menu-alignment-v1!}在三个已接触实例上做12局开发对照；随后十个已接触实例做40局扩展开发对照。新实现\passthrough{\lstinline!resource-scoped-response-v2!}使用四个新种子1042---1045，沿用Test\_00前四个配置重新生成保留实例，五方法各运行一次，共20局，清单见\href{../expr/datasets/flatland-core/v3-holdout/manifest.json}{保留数据manifest}。这些不是官方发布的原始排行榜测试文件。各批结果分别报告，不能合并成同版本独立样本，也不能把开发集改善当作保留收益。源码快照随run保存。上述批量实验均为规则响应、无LLM，不将软件Agent等同完整LLM
Agent；Qwen混合验证另见第5.4节。

\begin{longtable}[]{@{}ll@{}}
\toprule
方法 & 差别 \\
\midrule
\endhead
independent & 本地最小代价规划，无队友消息 \\
single & 同协议但每方只返回1项响应，组合冲突则拒绝 \\
menus\_no\_repair & 最多3项响应，不计修复代理成本 \\
menus & 最多3项响应，计1单位计划变更成本 \\
\bottomrule
\end{longtable}

同一路线候选基础、低层执行器和观测接口。独立方法每步可检查，协商方法受轮换发起机会限制，不能宣称决策机会完全相同。同一pickle复用不保证不同策略遇到完全相同的执行相关故障。

主指标为到达数和惩罚完成步数：每列车到达时刻相加，未完成列车按环境原生最大步数计。运行上限1000步，原生终止时结束。惩罚步数与到达数联合解释，少完成也可能因其他车更早而得到较低总和。另计消息数、UTF-8消息序列化字节、组合检查、墙钟和计划变更；字节不等于真实网络流量。

\hypertarget{ux5206ux5c42ux5b9eux9a8cux7ed3ux679c}{%
\chapter{分层实验结果}\label{ux5206ux5c42ux5b9eux9a8cux7ed3ux679c}}

\hypertarget{ux65e9ux671f12ux5c40ux5f00ux53d1ux5bf9ux7167ux7528ux4e8eux5f62ux6210ux8bbeux8ba1ux4e0dux4f5cux6700ux7ec8ux6cdbux5316ux7ed3ux8bba}{%
\section{早期12局开发对照：用于形成设计，不作最终泛化结论}\label{ux65e9ux671f12ux5c40ux5f00ux53d1ux5bf9ux7167ux7528ux4e8eux5f62ux6210ux8bbeux8ba1ux4e0dux4f5cux6700ux7ec8ux6cdbux5316ux7ed3ux8bba}}

\begin{longtable}[]{@{}llllll@{}}
\toprule
方法 & Level\_0：到达/惩罚步数 & Level\_1 & Level\_2 & 总到达 & 总惩罚步数 \\
\midrule
\endhead
independent & 3/7，1391 & 7/7，853 & 3/7，1139 & 13/21 & 3383 \\
single & 6/7，1086 & 7/7，923 & 7/7，923 & 20/21 & 2932 \\
menus\_no\_repair & 7/7，983 & 7/7，925 & 7/7，906 & 21/21 & 2814 \\
menus & 7/7，983 & 7/7，925 & 7/7，888 & 21/21 & 2796 \\
\bottomrule
\end{longtable}

多响应含修复成本相对独立规划，在这三个实例的惩罚步数降低约17.35\%；相对单响应降低约4.64\%。这是描述性开发集比较，不是显著性结论。Level\_1独立规划更快，协调有保守代价；修复权重在前两例无差别，只在第三例改善，因此不称所有模块各自都有效。

\hypertarget{ux65e9ux671fux5f00ux53d1ux5bf9ux7167ux7684ux6210ux672c}{%
\subsection{早期开发对照的成本}\label{ux65e9ux671fux5f00ux53d1ux5bf9ux7167ux7684ux6210ux672c}}

\begin{longtable}[]{@{}lllll@{}}
\toprule
方法 & 三局消息数 & 序列化消息字节 & 三局循环秒 & 已提交计划变更次数 \\
\midrule
\endhead
independent & 0 & 0 & 0.831 & 0（该指标分支未计，不代表无计划变化） \\
single & 4224 & 23312959 & 10.436 & 8 \\
menus\_no\_repair & 3355 & 21793898 & 12.294 & 16 \\
menus & 3285 & 21554064 & 11.978 & 14 \\
\bottomrule
\end{longtable}

菜单协调比单响应少部分消息，但比独立方法显著昂贵；大轨迹报文也不适合直接部署。墙钟含控制及相关日志但不是严格性能基准，不能将低于十几秒解释为模型推理速度。

\hypertarget{ux6269ux5c5540ux5c40ux5f00ux53d1ux5bf9ux7167}{%
\section{扩展40局开发对照}\label{ux6269ux5c5540ux5c40ux5f00ux53d1ux5bf9ux7167}}

十个开发实例、每例七列车，每方法运行十局。原始结果见\href{../expr/004-distributed-agent-core/runs/20261111-213049-aligned-ablation.json}{扩展批次}。这些实例均已接触，不是新增保留集，也不与早期12局相加计算样本量。

\begin{longtable}[]{@{}llllll@{}}
\toprule
方法 & 到达/总数 & 惩罚完成步数 & 消息字节 & 组合检查 & 循环秒 \\
\midrule
\endhead
independent & 33/70 & 11198 & 0 & 0 & 2.984 \\
single & 49/70 & 9954 & 116671092 & 4358 & 96.115 \\
menus\_no\_repair & 63/70 & 9200 & 91650580 & 83389 & 102.645 \\
menus & 63/70 & 9182 & 91410394 & 79431 & 108.017 \\
\bottomrule
\end{longtable}

这一批多响应比单响应多到达14列车；加入修复代价未改变到达数，仅使惩罚步数减少18。它支持继续验证响应菜单，但不能说明保留场景也有同样收益，计算开销仍明显。

\hypertarget{ux5c40ux65b0ux79cdux5b50ux4fddux7559ux5bf9ux7167ux4e3bux8981ux8bc4ux4ef7ux4f9dux636e}{%
\section{20局新种子保留对照：主要评价依据}\label{ux5c40ux65b0ux79cdux5b50ux4fddux7559ux5bf9ux7167ux4e3bux8981ux8bc4ux4ef7ux4f9dux636e}}

四个新种子在运行前固定，未根据本轮结果更换种子或调整参数。五方法使用相同场景文件并核对SHA256，全部失败保留。原始结果见\href{../expr/004-distributed-agent-core/runs/20261111-221706-aligned-ablation.json}{保留批次}，逐实例配对与限制见\href{../expr/004-distributed-agent-core/runs/20261111-221706-aligned-ablation.md}{批次报告}。

\begin{longtable}[]{@{}llllllll@{}}
\toprule
方法 & 到达/总数 & 惩罚完成步数 & 消息数 & 消息字节 & 组合检查 & 循环秒 &
已提交计划变更次数 \\
\midrule
\endhead
independent & 10/28 & 4634 & 0 & 0 & 0 & 0.824 & 未按同定义记录 \\
single & 15/28 & 4333 & 12938 & 84728639 & 2119 & 44.000 & 28 \\
menus\_no\_repair & 15/28 & 4371 & 12681 & 107478013 & 97880 & 67.116 & 58 \\
menus & 15/28 & 4371 & 12681 & 107477763 & 97301 & 67.898 & 58 \\
official\_dla & 28/28 & 3328 & 不适用 & 不适用 & 不适用 & 0.474 & 不适用 \\
\bottomrule
\end{longtable}

DLA不使用本协议，其日志中消息和修复为零不能解释成同通信接口下无成本。它是不同信息条件下的工程参照，不是协调模块的公平单因素消融。

保留结果改变了对开发结果的判断：单响应和两种多响应都比独立方法多到达5列车，但多响应没有比单响应多完成任务，惩罚步数反而增加38。完整多响应比单响应少257条消息，却多传输约22.75
MB序列化数据，组合检查由2119增至97301。因此``消息条数减少''不能直接写成``协调成本下降''。两种多响应到达数和惩罚步数完全相同，未体现修复代价的额外任务收益。

只有四个实例，每方法每例一次，列车之间也不是独立样本。本文不据28列车宣称统计显著性或稳定泛化。不同策略的执行相关故障仍可能不同，单次墙钟也不构成稳定性能基准。若后续据这些失败改方法，这批数据就应转为开发诊断，不能继续充当未接触的保留集。

\hypertarget{qwen2.5-1.5bux6df7ux5408ux9a8cux8bc1ux4e0dux80fdux8bc1ux660eux6a21ux578bux589eux76ca}{%
\section{Qwen2.5-1.5B混合验证：不能证明模型增益}\label{qwen2.5-1.5bux6df7ux5408ux9a8cux8bc1ux4e0dux80fdux8bc1ux660eux6a21ux578bux589eux76ca}}

在一个保留场景中，独立与完整多响应各运行一局，均使用qwen2.5:1.5b、温度0、最多6次真实调用，之后规则回退。结果见\href{../expr/004-distributed-agent-core/runs/qwen-holdout-pair.json}{Qwen配对记录}。

\begin{longtable}[]{@{}lllllll@{}}
\toprule
方法 & 到达 & 惩罚完成步数 & 模型调用 & 输入/输出token & 规则回退次数 & 墙钟秒 \\
\midrule
\endhead
independent & 3/7 & 958 & 6 & 2331 / 42 & 92 & 29.975 \\
menus & 7/7 & 800 & 6 & 2486 / 42 & 2669 & 25.012 \\
\bottomrule
\end{longtable}

该场景的规则方法取得相同到达数与惩罚步数，因此本次只说明小模型可以接入候选选择，不能证明LLM带来收益。两方法调用时机、输入候选及轨迹不同，虽然调用上限相同，仍不是严格同输入的框架开关消融；规则回退数也是不同运行过程的计数，不等于模型失败数。单次计时不能说明多响应稳定更快。

另有\href{../expr/004-distributed-agent-core/runs/qwen-matched-probe.json}{同输入局部探针}，使用三份相同局部输入检查模型选择，但没有改变端到端协调过程，不能作为完整方法消融。严格同信息、同调用规则的Qwen端到端对照仍属后续研究。

\hypertarget{ux6700ux7ec8ux9a8cux8bc1ux4fddux7559ux6709ux6548ux5b89ux6392ux6309ux9700ux534fux8c03ux7684ux964dux672cux4ef7ux503c}{%
\section{最终验证：保留有效安排、按需协调的降本价值}\label{ux6700ux7ec8ux9a8cux8bc1ux4fddux7559ux6709ux6548ux5b89ux6392ux6309ux9700ux534fux8c03ux7684ux964dux672cux4ef7ux503c}}

最终正面证据来自\href{../expr/009-route-permit/README.md}{009实验说明}，原始批次为\passthrough{\lstinline!009-route-permit/runs/20261114-154929-745791!}。复用三个已有Flatland七车场景，设置无、低、高三档固定时钟STOP扰动，四方法共36局；真实调用env.step，统一1000步上限并关闭原生故障。这是开发验证，不是新保留集，也不是原排行榜时限。STOP是动作指令，不保证运动中的车辆瞬间静止。

各车使用同一条固定路线；共享路线的车辆按编号持有互斥通行许可，不相交路线可并行。Agent分别答复是否允许入场，确认后许可持续有效。任务只修改permit片段，其他任务片段不重置。该方案较保守，最大同时在轨车辆数为2，未实现区段换序和复杂连锁修复。

四种方法为：全车串行；每8步询问许可；每步询问许可；等待且相关前驱完成时才按需询问。后三者使用相同许可规则、轮换发起及消息协议，差别是触发时机。所有方法均发布车辆状态变化，按接收者计消息和序列化字节，状态发布成本纳入总量；字节不是实际网络抓包或模型token。

\begin{longtable}[]{@{}lllll@{}}
\toprule
方法 & 九局总到达 & 累计完成步数 & 总通信字节（含状态） & 本地响应评估次数 \\
\midrule
\endhead
全车串行 & 63/63 & 20707 & 34704 & 0 \\
每8步询问 & 63/63 & 19163 & 1695256 & 2262 \\
每步询问 & 63/63 & 16277 & 10911441 & 14906 \\
按需协调 & 63/63 & 16277 & 379014 & 393 \\
\bottomrule
\end{longtable}

九组按需协调与每步询问的动作和位置轨迹逐步完全相同，见批次内\passthrough{\lstinline!trajectory-equivalence.json!}。在执行结果不变的情况下，通信减少\passthrough{\lstinline!1−379014/10911441约96.53\%!}，响应评估减少\passthrough{\lstinline!1−393/14906约97.36\%!}。这是本报告最明确的降本证据：不是少完成任务，也不是只少算了状态通知，而是减少了条件未变时的重复协调。

相对每8步询问，按需协调完成步数减少约15.06\%，通信减少约77.64\%；时间改善部分来自不必等待下次询问周期。相对全车串行，完成步数减少约21.39\%，但通信更多，因此不能说所有指标全面领先。响应评估次数是工作量代理，不等于相同比例的墙钟、FLOPs或token节省。所有运行均无LLM调用。

36局全部到达，轨迹同格/交换占用审计违规及未授权入场均为0；未涉及的other\_work片段版本保持不变。009的两项测试分别检查许可拒绝和一场景等轨迹降通信，另对完整九组保存的轨迹进行了等价核验。有限采样审计不构成形式化安全证明，63辆次也不是63个独立样本。

\textbf{结论限定为：在当前三个开发场景和同步可靠条件下，验证了保留有效安排、按需协调的降本价值。}
对照中周期重复询问较冗余，因此结果不能解释为超过成熟事件触发调度，更不能证明新颖性、LLM增益、多响应收益或完整框架有效。复杂连锁协商仍作为后续工作。

\hypertarget{ux771fux5b9eux534fux5546ux6848ux4f8bux4e0dux662fux53eaux589eux52a0ux672aux8c03ux7528ux63a5ux53e3}{%
\chapter{真实协商案例：不是只增加未调用接口}\label{ux771fux5b9eux534fux5546ux6848ux4f8bux4e0dux662fux53eaux589eux52a0ux672aux8c03ux7528ux63a5ux53e3}}

案例来自\href{../expr/004-distributed-agent-core/runs/20261110-145737-343494-flatland-distributed/summary.json}{Level\_0的menus运行}，提案键\passthrough{\lstinline!141:1:3!}，即第141环境步由列车1发起。按0起始编号，这些都是七列车中的不同Agent。

\begin{enumerate}
\def\labelenumi{\arabic{enumi}.}
\tightlist
\item
  发起方候选0公开32步占用、版本18和有效期至149；该候选的后续预测还包括159步以后的进入安排。\textbf{149后是预测，不是被该8步期限保证的通行承诺。}
  这是当前表示需要清楚区分的限制。
\item
  列车2首项响应路线经过(9,17)→(9,16)→(9,15)，而列车3反向从(9,15)向(9,16)移动，首选组合发生换位冲突。
\item
  列车2自己提供替代路线，经(9,17)→(8,17)→(8,16)。它不是列车1为列车2生成的路线。
\item
  \passthrough{\lstinline!negotiation.json!}记录候选0检查3个组合后找到相容项，选择列车2菜单编号2；列车2该项公开增量代价为0，前两项为-1，但前两项不能与队友形成相容组合。
\item
  同一提案记录selected\_combination与commit。动作日志显示第142步执行后列车2位于(8,17)，第143步位于(8,16)，替代路线实际进入执行；该车最终159步到达。
\end{enumerate}

此例选中的替代项修复代理成本为0，因而不能把这一次替代选择称为一次新增修复。它说明替代菜单、组合核验与提交执行确实联通，不证明没有替代项就必然任务失败，也不证明每次协商都更好。

\hypertarget{ux8865ux5145ux63a2ux7d22ux4e0eux5931ux8d25}{%
\chapter{补充探索与失败}\label{ux8865ux5145ux63a2ux7d22ux4e0eux5931ux8d25}}

\hypertarget{ux5386ux53f2qwenux63a5ux53e3ux8bd5ux8dd1}{%
\section{历史Qwen接口试跑}\label{ux5386ux53f2qwenux63a5ux53e3ux8bd5ux8dd1}}

早期一次单组混合验证有6次模型调用、2742输入token和42输出token，剩余1785次候选响应选择预算回退，6/7到达，约32.76秒。这是历史接口试跑，不与第5.4节后续配对验证混用，也不是新增的同条件消融；仅说明接口和成本记录可用。当前Qwen结论以第5.4节为准。

\hypertarget{ux5206ux7ec4ux5931ux8d25}{%
\section{分组失败}\label{ux5206ux7ec4ux5931ux8d25}}

按目的地分组不等于按协作关系分组。跨组共享轨道导致扩展与重启，加入轮换发起后仍出现7/7、7/7、1/7，消息1635、3128、4850。该失败提示：通信边界应基于交互内容与资源依赖，而不是为了省消息机械分组。当前不宣称分组节省通信。

\hypertarget{ux5b98ux65b9ux5de5ux7a0bux53c2ux7167}{%
\section{官方工程参照}\label{ux5b98ux65b9ux5de5ux7a0bux53c2ux7167}}

接入Flatland
Association的DeadLockAvoidancePolicy固定commit，在十个旧开发配置中DLA69/70、本方法63/70，惩罚步数8993对9182，循环约0.75秒对46.54秒。官方规则可使用全环境信息，是工程参照而非同信息消融，也不是已确认的SOTA。课程不要求胜过它，但这些负面证据说明当前设计尚不能替代成熟铁路调度器。

\hypertarget{ux5b66ux751fux63d0ux51faux7684ux8bbeux8ba1ux9009ux62e9ux53caux53cdux601d}{%
\chapter{学生提出的设计选择及反思}\label{ux5b66ux751fux63d0ux51faux7684ux8bbeux8ba1ux9009ux62e9ux53caux53cdux601d}}

以下依据实际讨论中的明确主张整理，不虚构访谈或事后实验预测；提交前仍需本人核阅。

\begin{itemize}
\tightlist
\item
  反对仅检查队友``是否已经开始''，要求覆盖未来依赖。
\item
  强调借鉴MPC内化为Agent能力，而不是模型调用控制器。
\item
  指出不应由A替B预测，而应让B对A候选给出自己的响应及代价。
\item
  要求至少三个成员、所有Agent同构核心、任务可插拔。
\item
  反对每次全网广播，提出任务组，再进一步按交互内容确定参与者。
\item
  指出课程交付不应被SOTA或专利目标无限扩大，要求回到具体课件点与完整报告。
\end{itemize}

AI协作也暴露了问题：早期把思想简化成保守拒绝；用接口存在代替能力完成；将不同版本结果并列；目的地分组缺少依据；把工程参照过度当成必要门槛。改进是保留原始失败、做同版本消融、记录可追踪案例，并明确``已实现''和``已证实''不同。

\hypertarget{ux7ed3ux8bbaux4e0eux540eux7eed}{%
\chapter{结论与后续}\label{ux7ed3ux8bbaux4e0eux540eux7eed}}

本作业从ReAct``再想''延伸到多Agent协作中的按需更新，实现了同构核心、相关方响应、条件确认及真实Flatland执行。最终009对照表明，已有可执行安排的外部条件未变时，不必重复协商：按需协调与每步询问九组执行轨迹完全相同，通信减少96.53\%，响应评估减少97.36\%。因此，本作业验证了保留有效安排、按需协调的降本价值，结论仅适用于所测开发场景、许可规则和同步可靠前提。

更广的框架目标尚未得到证明。004保留集两种多响应均15/28，与单响应相同且更昂贵，加入修复代价没有额外收益；Qwen验证不支持模型独立增益。005高冲突停滞、006在线短视失败促使我们补充回退和依赖保护，但恢复到达不等于证明创新优势。007/008只提供局部补丁及连锁协调的能力证据，009只验证全路线许可的按需更新，不能相互替代证据。

后续研究包括细粒度区段许可、连锁修复的真实执行、新种子保留测试及与成熟事件触发方法的比较。核心虽已提供分段补丁和其他扩展，具体任务是否正确接入必须逐项验证，不能把接口存在当作降本成立。这些不作为本次课程提交前置条件。

课程层面已经形成可解释设计、真实代码、对照结果、失败和反思；完整会话与Public链接仍按用户要求暂缓，因此本稿不是已满足全部硬性提交条件的最终发布版。论文级扩展不作为当前交付前置条件。

\hypertarget{ux590dux73b0ux4e0eux8bc1ux636eux7d22ux5f15}{%
\chapter{复现与证据索引}\label{ux590dux73b0ux4e0eux8bc1ux636eux7d22ux5f15}}

\begin{itemize}
\item
  最终正面验证：\href{../expr/009-route-permit/README.md}{009说明}；批次\passthrough{\lstinline!009-route-permit/runs/20261114-154929-745791/results.json!}；九组等轨迹审计为同目录\passthrough{\lstinline!trajectory-equivalence.json!}；源码及核心快照随批次保存。
\item
  过程与失败记录：\href{../expr/006-mpc-flatland/在线验收状态.md}{在线验收状态}、\href{../expr/006-mpc-flatland/滚动执行恢复验收.md}{滚动执行恢复}、\href{../expr/007-local-chain/README.md}{局部连锁能力}、\href{../expr/008-flatland-local-repair/README.md}{铁路离线修复}。
\item
  早期12局开发对照：\passthrough{\lstinline!004-distributed-agent-core/runs/20261110-145826-aligned-ablation.json!}。
\item
  扩展40局开发对照：\passthrough{\lstinline!004-distributed-agent-core/runs/20261111-213049-aligned-ablation.json!}。
\item
  主要保留验证：\passthrough{\lstinline!004-distributed-agent-core/runs/20261111-221706-aligned-ablation.json!}；过程与不确定性记录为同目录\passthrough{\lstinline!20261111-221706-aligned-ablation-audit.md!}。
\item
  Qwen两局验证与局部探针：\passthrough{\lstinline!004-distributed-agent-core/runs/qwen-holdout-pair.json!}、\passthrough{\lstinline!004-distributed-agent-core/runs/qwen-matched-probe.json!}。
\item
  每run保留\passthrough{\lstinline!summary.json!}、\passthrough{\lstinline!negotiation.json!}、\passthrough{\lstinline!messages.json!}、\passthrough{\lstinline!trajectory.jsonl!}及源码快照；历史复核应读取对应版本，而不是直接用当前代码替代。
\item
  批量入口：\passthrough{\lstinline!004-distributed-agent-core/scripts/run\_aligned\_ablation.py!}；任务适配器：\passthrough{\lstinline!004-distributed-agent-core/src/flatland\_task.py!}。
\item
  官方参照：\passthrough{\lstinline!004-distributed-agent-core/官方强基线十场景对照.md!}。
\item
  组与模型验证：\passthrough{\lstinline!004-distributed-agent-core/Qwen与任务分组验证.md!}。
\item
  前提：\href{../expr/前提.md}{前提.md}。无密钥写入报告，无需付费API。
\end{itemize}

\begin{lstlisting}
$py = "最终作业/expr/003-flatland-core/.venv/Scripts/python.exe"
& $py -m unittest discover -s 最终作业/expr/004-distributed-agent-core/tests -v
& $py 最终作业/expr/004-distributed-agent-core/scripts/run_aligned_ablation.py
& $py -m unittest discover -s 最终作业/expr/009-route-permit -v
& $py 最终作业/expr/009-route-permit/run.py
\end{lstlisting}

当前代码允许更多对照；复核历史主实验时以各run源码快照为准。重新运行会创建新目录，不覆盖旧证据。

历史复核：2026-11-10，17项测试通过，Level\_0菜单方法复跑7/7、惩罚完成步数983，与早期开发表一致；运行目录\passthrough{\lstinline!004-distributed-agent-core/runs/20261110-155712-922411-flatland-distributed/!}，循环约4.59秒。后续v2验证记录为18项测试通过，见本轮验证结论。这些是当时的测试记录，本次文档更新没有重新运行实验，也不把代表实例复跑当作新增独立样本。

\hypertarget{ux53c2ux8003ux8d44ux6599ux4e0eaiux4f7fux7528ux8bf4ux660e}{%
\chapter{参考资料与AI使用说明}\label{ux53c2ux8003ux8d44ux6599ux4e0eaiux4f7fux7528ux8bf4ux660e}}

\begin{enumerate}
\def\labelenumi{\arabic{enumi}.}
\tightlist
\item
  课程课件，第7讲课时1，第8页ReAct节拍；第4页``模型出主意，应用来动手''为辅助边界。
\item
  Yao et al.~ReAct: Synergizing Reasoning and Acting in Language Models，ICLR
  2023，https://arxiv.org/abs/2210.03629；官方代码https://github.com/ysymyth/ReAct。
\item
  Rawlings, Mayne, Diehl. Model Predictive Control: Theory, Computation, and Design.
  2nd ed., 2017。作为有限前瞻与滚动控制背景，不意味着本实现具备其理论保证。
\item
  Flatland环境：https://github.com/flatland-association/flatland-rl；场景：https://github.com/flatland-association/flatland-scenarios；工程基线：https://github.com/flatland-association/flatland-baselines。
\item
  其他多Agent计划修复与协调阅读见\href{../expr/../参考文献/多agent协作/总览.md}{文献总览}，不是本报告新颖性已获证明的依据。
\end{enumerate}

本报告由用户与AI多轮讨论、实现和核验形成，AI参与代码与文字整理；用户的关键设计判断见第8节。现有会话文件不被声称完整；按用户指示，本轮不导出、不补造、不上传会话，不创建或设置公开仓库。
